方程要连同定解条件一起写.分离变量的系数不算完,就不能算这道题做完.

\begin{problemset}
\renewcommand{\theenumi}{\thechapter.\arabic{enumi}}
\renewcommand{\labelenumi}{\theenumi.}
\item 由牛顿定律写出两端张力为 $T$、线密度为 $\rho$ 的横振动方程,并指出波速 $a$ 与 $T,\rho$ 的关系.说明固定端和自由端分别对应哪一类边界条件.

\item 判断下列方程的类型,并说明哪一类数据是适当的: $u_{tt}-a^2 u_{xx}=0$, $u_t-a^2 u_{xx}=0$, $u_{xx}+u_{yy}=0$.

\item 无界弦 $u_{tt}-a^2 u_{xx}=0$, $u(x,0)=\sin(kx)$, $u_t(x,0)=0$.用达朗贝尔公式写出 $u(x,t)$,并验证它满足方程和初值.

\item 两端固定的弦, $u(0,t)=u(\ell,t)=0$, $u(x,0)=\varphi(x)$, $u_t(x,0)=0$.分离变量,写出本征值和级数解.取 $\varphi(x)=x(\ell-x)$,把系数积分出来.

\item 第三类边界 $X(0)=0$, $X'(l)+hX(l)=0$, $h>0$.证明本征值由 $\tan(\mu l)=-\mu/h$ 决定,并且每个区间 $((n-1/2)\pi,n\pi)$ 里恰有一根.

\item 无限长杆 $u_t=a^2 u_{xx}$, $u(x,0)=\varphi(x)$.写出热核卷积.若 $\varphi$ 是常数 $T_0$,说明 $u$ 为什么保持为 $T_0$.

\item 在 $0<x<\ell$ 上求 $X''+\lambda X=0$ 的本征值和本征函数:两端自由;左端固定、右端自由.说明第二种情形为什么没有 $\lambda=0$.

\item 两端第三类条件 $X'(0)=\alpha X(0)$, $X'(\ell)=-\beta X(\ell)$, $\alpha,\beta\geq 0$ 且不同时为零.由能量恒等式说明 $\lambda>0$,并推出 $\tan(\omega\ell)=(\alpha+\beta)\omega/(\omega^2-\alpha\beta)$.再说明 $\alpha=\beta$ 且 $\alpha\ell<-2$ 时为什么有两个负本征值.

\item 两端夹紧的梁,$X(0)=X'(0)=X(\ell)=X'(\ell)=0$.把通解里的两个系数用左端条件消去,证明频率满足 $\cosh(k\ell)\cos(k\ell)=1$.

\item $u_t=u_{xx}$, $0<x<1$, $u(0,t)=0$, $u(1,t)=1$, $u(x,0)=0$.先求稳态,再把差按正弦级数展开.

\item 矩形 $0<x<a$, $0<y<b$ 上 $\Delta u=0$,左右和下边为零,上边是 $\sin(\pi x/a)$.写出解,并在四条边上核对.

\item 矩形膜 $0<x<a$, $0<y<b$,四周固定.写出 $u_{tt}=c^2\Delta u$ 的频率 $\omega_{mn}$.说明 $a=b$ 时哪些模式简并.

\item $u_{tt}-2a u_{tx}-u_{xx}=0$, $0<x<\pi$, $u(0,t)=u(\pi,t)=0$.用 $e^{\imath k(t-ax)}\sin(nx)$ 定出 $k$,并说明为什么普通的 $X(x)T(t)$ 乘积解不出这个方程.
\end{problemset}
